\section{What is open, stated exactly}

Erdős \#249 asks whether $S = \sum_{n\ge 0} \varphi(n)/2^n$ is irrational. It is
\ev{Open}: nothing in this corpus decides it.

Every attack recorded in Parts
I--IV terminates, in one of exactly two ways, at a small number of
\emph{cofinal supply obligations} on the totient tail
$R_N := \sum_{j\ge 0}\varphi(N+1+j)/2^{j+1}$
(\lean{Erdos249257.TotientTailPeriodKiller.totientTail}{TotientTailPeriodKiller.lean:57}):
either it proves a genuine theorem of the shape ``\emph{if} this cofinal
predicate holds, \emph{then} $S$ is irrational,'' with the predicate itself
left completely unproved, or it prunes one specific proof shape and shows it
cannot supply the predicate.

This section renders every such obligation
exactly as Lean states it, gives its site, records which of them are proved
\emph{equivalent} to one another (iff) and which are proved only
\emph{sufficient} (one-directional, an easier target implied by a harder
one), and closes by naming, without overclaiming, the piece of new
mathematics each surviving form would need. Quantifier order is preserved
exactly as written in Lean; nowhere below does a finite verified list stand
in for a cofinal supply, and nowhere does an implication get reported as a
solved case.

Throughout, $H_t := \mathrm{periodLcm}(t) = \operatorname{lcm}(1,\dots,t)$
(\lean{Erdos249257.CarrySurvivorExtinction.periodLcm}{CarrySurvivorExtinction.lean:515}),
$D(h,N,L) := \sum_{j<L}\bigl(\varphi(N{+}h{+}1{+}j)-\varphi(N{+}1{+}j)\bigr)\cdot 2^{L-1-j}\in\mathbb{Z}$
is \texttt{windowDiscrepancy}
(\lean{Erdos249257.TotientTailPeriodKiller.windowDiscrepancy}{TotientTailPeriodKiller.lean:66}),
and the base predicate named ``$\mathrm{Sep}(h,N,L)$'' in the task brief is
literally \texttt{certifiedKill}:
\[
  \mathrm{certifiedKill}(h,N,L) :\equiv
  (N+h+L+2 : \mathbb{Z}) < D(h,N,L) \bmod 2^L
  \ \wedge\
  D(h,N,L) \bmod 2^L < 2^L - (N+h+L+2).
\]
\lean{Erdos249257.TotientTailPeriodKiller.certifiedKill}{TotientTailPeriodKiller.lean:72}
\quad \ev{Lean} (the definition; decidable, instance at line 84)
\quad \scale{n/a} \quad \coord{other:binary-window}. In words: the residue
of the window discrepancy modulo $2^L$ avoids a shrinking radius-$(N+h+L+2)$
neighbourhood of $0$ inside the growing modulus $2^L$.

\subsection{The certificate-completeness converter}

Every reformulation below ultimately measures the same underlying real
quantity, because certificates are \emph{complete} receipts of
non-integrality, not merely sufficient ones:
\[
  \bigl(\exists L,\ \mathrm{certifiedKill}(h,N,L)\bigr)
  \iff
  R_{N+h} - R_N \notin \operatorname{range}\bigl((\uparrow)\colon \mathbb{Z}\to\mathbb{R}\bigr).
\]
\lean{Erdos249257.exists_certifiedKill_iff_tail_diff_notMem_int}{LcmConeFlatness.lean:316}
(wrapped as \texttt{totient\_tail\_window\_kill\_exists\_iff\_tail\_diff\allowbreak\_not\_int} at
\lean{Erdos249257.totient\_tail\_window\_kill\_exists\_iff\_tail\_diff\_not\_int}{CertificateKernel.lean:19025})
\quad \ev{Lean} \quad \scale{uniform} (holds for all $h,N$)
\quad \coord{other:binary-window}.

\emph{Note.} This iff is the reason every ``certificate supply'' form below
can be restated, without loss, in ``pure non-integrality'' language with no
certificate vocabulary at all (the B8/TE forms). It does \emph{not} assert
that the cofinal supply exists; it only says the certificate route and the
real-analytic route are the same route.

\subsection{The base cofinal obligation — the wall}

\begin{quote}
\textbf{Obligation 1 (certificate-supply normal form).}
\[
  \mathrm{Irrational}(S)
  \quad\Longleftrightarrow\quad
  \forall h\ge 1,\ \forall N_0,\ \exists N\ge N_0,\ \exists L,\
  \mathrm{certifiedKill}(h,N,L).
\]
\end{quote}
This is exactly the quantifier structure named ``$\mathrm{Sep}(h,N,L)$'' in
the task brief: $\forall h\ge 1\ \forall N_0\ge 0\ \exists N\ge N_0\ \exists L\ \mathrm{Sep}(h,N,L)$.
\lean{Erdos249257.irrational_totient_series_iff_certificate_supply}{LcmConeFlatness.lean:412}
\quad \ev{Lean} (the equivalence is proved; the supply itself is unsupplied)
\quad \scale{cofinal} \quad \coord{other:binary-window}.

The reverse direction is by contradiction against the tail-period law: if $S$ were
rational, \texttt{eventual\_period\_of\_not\_irrational}
(\lean{Erdos249257.eventual_period_of_not_irrational}{TotientTailPeriodKiller.lean:358},
\ev{Lean}, unconditional) supplies a period $h>0$ and pre-period $N_0$ with
$R_{N+h}-R_N\in\mathbb{Z}$ for all $N\ge N_0$; the certificate hypothesis
instantiated at that $(h,N_0)$ produces an $N\ge N_0$ where
$R_{N+h}-R_N\notin\mathbb{Z}$, a contradiction via \texttt{tail\_diff\_notMem\_int\_of\_certifiedKill}
(\lean{Erdos249257.tail_diff_notMem_int_of_certifiedKill}{TotientTailPeriodKiller.lean:262}).

For the forward direction, irrationality gives pointwise non-integrality of
every positive tail shift and certificate completeness supplies a witness
already at $N=N_0$.

\emph{What supplying Obligation 1 would take:} an arithmetic (not merely
existential) argument that, for every period length $h$ and past every
threshold $N_0$, the totient window discrepancy's residue mod $2^L$ can be
pushed off $0$ by a margin growing with the modulus — i.e., genuine
cancellation/anti-concentration information about $\varphi$ on a sliding
window, at every scale, for every $h$. No such argument exists in the
corpus; the rows below are exact normal forms or sufficient producers for
this missing input, never a proof that the input holds.

\subsection{Collapsing free parameters: the multiple / diagonal / cone chain}

Three further reformulations replace Obligation 1's two free parameters
$(h,N_0)$ with successively less information.  The multiple and cone
predicates have explicit reductions to irrationality; the lcm-diagonal
predicate has a registered iff.  Since Obligation 1 itself is an iff, every
predicate both implied by Obligation 1 and sufficient for irrationality is
also propositionally equivalent to it, even when the corpus exposes only
the two directed constructions rather than a separately named iff.

\paragraph{Multiple-period collapse (formally looser, but endpoint-equivalent).}
\[
  \forall h_0>0,\ \forall N_0,\ \exists m>0,\ \exists N\ge N_0,\ \exists L,\
    \mathrm{certifiedKill}(m\cdot h_0, N, L)
  \ \Longrightarrow\ \mathrm{Irrational}(S).
\]
\lean{Erdos249257.irrational_totient_series_of_multiple_certificate_supply}{CarrySurvivorExtinction.lean:502}
\quad \ev{Lean} \quad \scale{cofinal} \quad \coord{other:lcm-period-multiple}.
Obligation 1's hypothesis trivially implies this one (take $m=1$), while the
displayed theorem sends it back to irrationality and hence, by the base iff,
back to Obligation 1.  It is an easier-looking target, not progress; it
remains exactly as open as \#249.

\paragraph{Diagonal collapse — one free parameter.}
\[
  \mathrm{Irrational}(S)
  \ \Longleftrightarrow\
  \forall t_0,\ \exists t\ge t_0,\ \exists L,\
    \mathrm{certifiedKill}(H_t, H_t, L).
\]
\lean{Erdos249257.irrational_totient_series_iff_lcm_diagonal_certificate_supply}{LcmConeFlatness.lean:426}
\quad \ev{Lean} \quad \scale{cofinal} \quad \coord{other:lcm-diagonal}.
This is the canonical single-quantifier restatement: write
$P(t) :\equiv \exists L,\ \mathrm{certifiedKill}(H_t,H_t,L)$; the obligation
is $\forall t_0\,\exists t\ge t_0,\ P(t)$. The reduction is a genuine proof,
not a relabelling: given $t\ge\max(h_0,N_0)$, $h_0\mid H_t$ and $H_t\ge
t\ge N_0$ simultaneously, so a single diagonal witness at $t$ discharges
\emph{both} of Obligation 1's free parameters at once, via the intermediate
lemma \lean{Erdos249257.irrational_totient_series_of_lcm_certificate_supply}{LcmDiagonalReduction.lean:92}
which itself reduces to the multiple-period form above.  Conversely,
irrationality and pointwise certificate completeness give the diagonal
witness already at $t=t_0$.

$P(t)$ is verified for every $t\le82$
(\lean{ErdosProblems.Skip.LadderT67.exists_diagonalKill_le_82}{ErdosProblems/Skip/LadderT67.lean:71264}).
The historical bank contained 28 explicit values through $t=64$
(\lean{Erdos249257.certifiedKill_diagonal_all_imported}{DiagonalPincerCertificates.lean:2870},
depths $[6,5,7,7,9,14,15,14,21,22,23,26,\dots]$ for
$t\in\{1,2,3,4,5,7,8,9,11,13,16,17,\dots\}$; endpoint
\lean{Erdos249257.certifiedKill_diagonal_t64}{DiagonalPincerCertificateT64Endpoint.lean:1928})
\ev{Cert} \scale{fixed} \coord{other:lcm-diagonal}).  The current contiguous
band is still a finite floor, not a cofinal supply; no certificate at $t=83$
is claimed.

\paragraph{Cone collapse — two multipliers, still the same wall.}
\[
  \forall t_0,\ \exists t\ge t_0,\ \exists q,m,L,\ 0<q\ \wedge\
    \mathrm{certifiedKill}(m\cdot H_t,\, q\cdot H_t,\, L)
  \ \Longrightarrow\ \mathrm{Irrational}(S).
\]
\lean{Erdos249257.irrational_totient_series_of_lcm_cone_window_kill_supply}{CertificateKernel.lean:19014}
\quad \ev{Lean} \quad \scale{cofinal} \quad \coord{other:lcm-cone}. The
diagonal form above is exactly the cell $q=m=1$, so any diagonal witness
trivially witnesses the cone form.  The cone predicate is therefore a
formally looser target, but its sufficiency theorem and the base iff make it
propositionally equivalent to \#249 rather than an unconditional advance. It
rests on a genuine
strengthening of the tail-period law, \emph{lcm-cone flatness}: if $S$ is
rational there is $t_1$ such that for every $t\ge t_1$ and every $q>0,m$,
$R_{qH_t+mH_t}-R_{qH_t}\in\mathbb{Z}$ — rationality flattens the
\emph{whole} cone $\{k H_t : k\ge 1\}$, not just one difference
(\lean{Erdos249257.rational_totient_series_forces_lcm_cone_flatness}{CertificateKernel.lean:19002},
\ev{Lean}, unconditional).

A sharper cone-form producer,
\texttt{coneNonflatCert}, needs only a \emph{one-sided} radius per vertex —
information-theoretically half of \texttt{certifiedKill}'s pairwise floor —
and proves that some pair in a finite multiplier menu $Q$ is non-integral;
the corresponding supply,
$\exists$ an unbounded-scale menu $Q$ with \texttt{coneNonflatCert} firing,
is \ev{Open}, \scale{cofinal},
\lean{Erdos249257.irrational_totient_series_of_lcm_cone_nonflat_supply}{CertificateKernel.lean:19140}.

\paragraph{Pure non-integrality form — certificate vocabulary stripped out.}
Via the completeness iff above, the diagonal and cone forms restate,
\emph{exactly}, with no reference to $\mathrm{certifiedKill}$ at all:
\[
  \forall t_0,\ \exists t\ge t_0,\ R_{2H_t}-R_{H_t}\notin\mathbb{Z}
  \ \Longleftrightarrow\ \text{diagonal obligation above}
  \ \Longrightarrow\ \mathrm{Irrational}(S),
\]
\lean{Erdos249257.irrational_totient_series_of_lcm_diagonal_nonintegrality_supply}{CertificateKernel.lean:19034}
(cone analogue at
\lean{Erdos249257.irrational\_totient\_series\_of\_lcm\_cone\_nonintegrality\_supply}{CertificateKernel.lean:19046})
\quad \ev{Lean} \quad \scale{cofinal} \quad \coord{other:real-analytic-nonintegrality}.
This is the frontier of \#249 with every piece of certificate machinery
removed: \emph{does $R_{2H_t}-R_{H_t}\notin\mathbb{Z}$ for infinitely many
$t$?} Nothing decides this either.

\subsection{The exponential-sum form}

A first-harmonic (Weyl-sum) cancellation statement is proved sufficient for
Obligation 1 by an elementary pigeonhole argument, with no case analysis and
no scale-degrading constants:
\[
  \mathrm{DTWFirstHarmonicNormGap} :\equiv\
  \forall h>0,\ \forall X_0,\ \exists X,L,\quad
  \max(X_0,1)\le X\ \wedge\
  16\,(2X+h+L+2)\le 2^L\ \wedge\
  \Bigl\|\sum_{N\in[X,2X)} e\bigl(D(h,N,L)/2^L\bigr)\Bigr\| \le \tfrac{21}{25}X.
\]
\lean{Erdos249257.DTWFirstHarmonicNormGap}{FirstHarmonicPivot.lean:75}
\quad \ev{Open} (the predicate; unproved at every $h,X_0$)
\quad \scale{cofinal} \quad \coord{other:first-harmonic}.

\[
  \mathrm{DTWFirstHarmonicNormGap} \ \Longrightarrow\ \mathrm{Irrational}(S).
\]
\lean{Erdos249257.irrational_totient_series_of_first_harmonic_norm_gap}{FirstHarmonicPivot.lean:83}
\quad \ev{Lean} \quad \scale{cofinal} \quad \coord{other:first-harmonic}. The
consumer already has Obligation 1's exact free-parameter shape ($X$ is a
completely free threshold, so applying the gap at $X\ge N_0$ gives
$\exists N\ge N_0\ \exists L$ directly); what is missing is not scale and not
coordinate, it is the arithmetic input itself — not one instance of a
constant-saving cancellation bound for the complex exponential sum
$\sum_{N\in[X,2X)} e(D(h,N,L)/2^L)$ is proved anywhere in the corpus, at any
$h,X,L$.

A strictly stronger \emph{subset} form (a saving on any nonempty
finite $T\subseteq[0,2X)$, no density or partition hypothesis) is also
proved sufficient and unconditional as an implication:
\lean{Erdos249257.exists_certifiedKill_of_first_harmonic_norm_gap}{FirstHarmonicPivot.lean:53}
consumes the elementary block lemma
\lean{Erdos249257.exists_certifiedKill_of_first_harmonic_gap}{FirstHarmonicGap.lean:128}
(\ev{Lean}, unconditional: \emph{any} constant-saving first-harmonic gap on
one dyadic block forces a finite kill certificate, via $\cos(\pi/8)>9/10$
and averaging).

\emph{What supplying this form would take:} genuine
cancellation for the first additive character of the totient window
discrepancy — a Weyl-type bound, uniform in $h$ and growing $L$, on a
character sum built from $\varphi$ differences, not merely an existence
statement. \lean{Erdos249257.ResidualGaugeObstruction.det_residualMonomialMatrix}{ResidualGaugeObstruction.lean:47}
rules out one entire class of candidate proofs: a residual-blind
determinant/rank certificate cannot see this cancellation, because an
explicit ``locked'' unit-norm gauge reconstructs the exact bad configuration
while keeping any Vandermonde-shaped minor nonzero — any real argument here
must couple the rows arithmetically, not merely normalise columns.

\subsection{The actual-orbit reformulation and the top-edge staircase}

A second, independently developed lane (the public pinned modules
\texttt{TotientActual*}/\texttt{TotientFixedRank*}) restates \#249 in terms of the \emph{actual}
power-of-two LCM diagonal, $H = H_{2^a}$, and proves a genuine \textbf{iff}
with the totient series itself — the strongest form of ``exact statement of
what remains'' anywhere in the corpus.

\begin{quote}
\textbf{The exact equivalence.}
\[
  \mathrm{Irrational}(S)
  \iff
  \forall a_0,\ \exists a\ge a_0,\ \mathrm{actualLcmTailOrbit}(a)\notin\mathbb{Z}.
\]
\end{quote}
\lean{Erdos257PeriodNoncollapse.DiagonalFreshLossBridge.PowerTwoOddWindowAffine.irrational_totientSeries_iff_actualLcmOrbitNonintegralitySupply}{TotientActualLcmOrbitNonintegrality.lean:37}
\quad \ev{Lean} \quad \scale{n/a} (an unconditional equivalence, not a
supply) \quad \coord{mobius-mersenne}, where
$\mathrm{actualLcmTailOrbit}(a) := 2^H(2^H-1)S - (\mathrm{totientPrefix}(2H)-\mathrm{totientPrefix}(H))$,
$H=H_{2^a}$
(\lean{Erdos257PeriodNoncollapse.DiagonalFreshLossBridge.PowerTwoOddWindowAffine.actualLcmTailOrbit_eq_scaled_totientSeries_sub_prefix}{TotientActualLcmOrbitSeparation.lean:68}).
\emph{This is \#249 restated with nothing left over}: no auxiliary
hypothesis, no scale caveat — cofinal non-integrality of this one sparse
sequence at power-of-two heights is exactly Erdős \#249. Everything else in
this subsection is an attempt to reach the right-hand side.

\paragraph{Route-pruning: total staircase collapse is impossible.}
Before any positive target, the file proves one entire proof shape is
empty. For $a\ge 8$ and room bound $J+K+(a{+}6)<2\cdot 2^a$ with a wide
enough modulus $2H+J+K+2<2^m$, a terminal window where \emph{every} one of
the last $m$ letters vanishes mod its own growing power of two is
impossible:
\[
  \neg\,\mathrm{ActualLcmTerminalDyadicStaircase}(a,J,K,m).
\]
\lean{Erdos257PeriodNoncollapse.DiagonalFreshLossBridge.PowerTwoOddWindowAffine.not_actualLcmTerminalDyadicStaircase_of_room}{TotientActualLcmTopEdgeStaircase.lean:251}
\quad \ev{Lean} \quad \scale{bounded} \quad \coord{mobius-mersenne}. The
mechanism is a bare positivity/divisibility contradiction (a positive
quantity below a modulus, divisible by that modulus, must be $0$) against
the unconditional positivity of every short-window letter,
$0<\mathrm{lcmRayArithmeticLetter}(2^a,j)$ for $a\ge8$, $0<j<2\cdot2^a$
(\lean{Erdos257PeriodNoncollapse.DiagonalFreshLossBridge.PowerTwoOddWindowAffine.lcmRayArithmeticLetter_pos_of_lt_two_mul}{TotientActualLcmOrbitArithmetic.lean:1703},
\ev{Lean}, unconditional, no rationality hypothesis).

\emph{Do not attempt a
full terminal-staircase producer}: it is provably empty. The corpus's own
route past this dead end is the \emph{punctured} staircase (all but the
last letter vanish; the last is retained and pinned exactly to the half-turn
$2^{m-1}$,
\lean{Erdos257PeriodNoncollapse.DiagonalFreshLossBridge.PowerTwoOddWindowAffine.puncturedDyadicStaircase_penultimate_eq_half}{TotientActualLcmTopEdgeStaircase.lean:1187},
\ev{Lean}, \scale{bounded}) and the one-sided residue-gap family below.

\paragraph{The surviving one-sided target.}
\[
  \mathrm{ActualLcmTopEdgeResidueGap}(a,J,K,m) :\equiv
  m\le K\ \wedge\
  2H+J+K+2 < 2^m\ \wedge\
  D(H,\,H+J,\,K) \bmod 2^m \le 2^m - (2H+J+K+2).
\]
\lean{Erdos257PeriodNoncollapse.DiagonalFreshLossBridge.PowerTwoOddWindowAffine.ActualLcmTopEdgeResidueGap}{TotientActualLcmTopEdgeStaircase.lean:900}
\quad \ev{Open} (a definition, unproved at any $a$)
\quad \scale{bounded} \quad \coord{mobius-mersenne}.

Note this is a
\emph{one-sided} inequality — only the upper (positive) carry arc is
excluded, not a symmetric two-sided band — because the corridor's lower half
is already discharged unconditionally: for $a\ge8$ and $J+(a{+}6)<2\cdot2^a$,
\[
  0 < R_{2H+J}-R_{H+J}
\]
\lean{Erdos257PeriodNoncollapse.DiagonalFreshLossBridge.PowerTwoOddWindowAffine.actualLcmTailDiff_shift_pos}{TotientActualLcmOrbitSign.lean:39}
\quad \ev{Lean} \quad \scale{bounded} \quad \coord{mobius-mersenne} — a rare
genuinely unconditional, non-hypothetical real-analytic theorem in this
corpus, needing no rationality assumption at all.

Its companion shows that
\emph{if} the orbit is integral, the residue is forced to the exact top
edge $2^K-e$ ($e>0$ small), \emph{outside} the arc a symmetric certificate
would need:
\lean{Erdos257PeriodNoncollapse.DiagonalFreshLossBridge.PowerTwoOddWindowAffine.actualLcm_integral_forces_topEdgeResidue}{TotientActualLcmOrbitSign.lean:211}
(\ev{Lean}, \scale{bounded}) — this is the exact statement of ``here is
what remains,'' pinning the missing exclusion to one named residue class
rather than leaving it implicit.

\[
  \mathrm{ActualLcmTopEdgeResidueGap}(a,J,K,m)\ \Longrightarrow\
  R_{2H+J}-R_{H+J}\notin\mathbb{Z}.
\]
\lean{Erdos257PeriodNoncollapse.DiagonalFreshLossBridge.PowerTwoOddWindowAffine.actualLcmTailDiff_notMem_int_of_topEdgeResidueGap}{TotientActualLcmTopEdgeStaircase.lean:1260}
\quad \ev{Lean} \quad \scale{bounded} \quad \coord{mobius-mersenne}.

The
cofinal target built from it:
\[
  \mathrm{PowerTwoActualLcmTopEdgeResidueGapSupply} :\equiv\
  \forall a_0,\ \exists a,K,m,\quad
  a_0\le a\ \wedge\ 8\le a\ \wedge\ K+(a{+}6)<2\cdot 2^a\ \wedge\
  \mathrm{ActualLcmTopEdgeResidueGap}(a,0,K,m).
\]
\lean{Erdos257PeriodNoncollapse.DiagonalFreshLossBridge.PowerTwoOddWindowAffine.PowerTwoActualLcmTopEdgeResidueGapSupply}{TotientActualLcmTopEdgeStaircase.lean:1325}
\quad \ev{Open} \quad \scale{cofinal} \quad \coord{mobius-mersenne}.
\[
  \mathrm{PowerTwoActualLcmTopEdgeResidueGapSupply} \Longrightarrow \mathrm{Irrational}(S).
\]
\lean{Erdos257PeriodNoncollapse.DiagonalFreshLossBridge.PowerTwoOddWindowAffine.irrational_totientSeries_of_topEdgeResidueGapSupply}{TotientActualLcmTopEdgeStaircase.lean:2184}
\quad \ev{Lean} \quad \scale{cofinal} \quad \coord{mobius-mersenne}.

\paragraph{Five strictly weaker links, each independently proved sufficient.}
The same file proves five further cofinal predicates, each strictly weaker
than \texttt{PowerTwoActualLcmTopEdgeResidueGapSupply} (each implies it, so
each is an easier target), with a direct \#249 endpoint of its own —
proving the \emph{weakest} of the six closes the entire cluster.

\begin{enumerate}
\item[(i)] $\mathrm{PowerTwoAdjacentSuffixMidbandSupply}$: replaces the
$m$-bit residue test with a symmetric two-sided band on the
\emph{adjacent-suffix} residue at depth $m$, buffered by the larger depth
$m{+}1$ so either branch stays inside the sign corridor. \lean{Erdos257PeriodNoncollapse.DiagonalFreshLossBridge.PowerTwoOddWindowAffine.PowerTwoAdjacentSuffixMidbandSupply}{TotientActualLcmTopEdgeStaircase.lean:1334}
\ev{Open} \scale{cofinal} \coord{mobius-mersenne}. Sufficiency:
\lean{Erdos257PeriodNoncollapse.DiagonalFreshLossBridge.PowerTwoOddWindowAffine.powerTwoActualLcmTopEdgeResidueGapSupply_of_adjacentSuffixMidband}{TotientActualLcmTopEdgeStaircase.lean:2108}, direct endpoint
\lean{Erdos257PeriodNoncollapse.DiagonalFreshLossBridge.PowerTwoOddWindowAffine.irrational_totientSeries_of_adjacentSuffixMidbandSupply}{TotientActualLcmTopEdgeStaircase.lean:2191}.

\item[(ii)] $\mathrm{PowerTwoOddGuardTopEdgeHalfWordBandSupply}$: at odd
depth $m=2q+1$ the adjacent-suffix residue is exactly twice a half-word
residue, and both directed edge widths halve to the same threshold
$H+q+2$:
\[
  H+q+2 \le \bigl(\text{half-word correction word}\bigr) \bmod 4^q
  \le 4^q - (H+q+2).
\]
\lean{Erdos257PeriodNoncollapse.DiagonalFreshLossBridge.PowerTwoOddWindowAffine.PowerTwoOddGuardTopEdgeHalfWordBandSupply}{TotientActualLcmTopEdgeStaircase.lean:1349}
\ev{Open} \scale{cofinal} \coord{mobius-mersenne}. This is a substantially
weaker demand than the earlier fixed $1/32$ central band elsewhere in the
corpus. Sufficiency:
\lean{Erdos257PeriodNoncollapse.DiagonalFreshLossBridge.PowerTwoOddWindowAffine.powerTwoAdjacentSuffixMidbandSupply_of_oddGuardTopEdgeHalfWordBand}{TotientActualLcmTopEdgeStaircase.lean:2057}.

\item[(iii)] $\mathrm{PowerTwoActualFinalTopEdgeMagnitudeSupply}$: the exact
centered-lift restatement of (ii),
\[
  H+q+2 \le |\mathrm{actualOddHalfCenteredLift}(a,q)|,
\]
\lean{Erdos257PeriodNoncollapse.DiagonalFreshLossBridge.PowerTwoOddWindowAffine.PowerTwoActualFinalTopEdgeMagnitudeSupply}{TotientActualLcmTopEdgeStaircase.lean:1471}
\ev{Open} \scale{cofinal} \coord{mobius-mersenne}. Proved
\textbf{equivalent} (an iff, not merely sufficient) to (ii):
\[
  \mathrm{PowerTwoOddGuardTopEdgeHalfWordBandSupply}
  \iff
  \mathrm{PowerTwoActualFinalTopEdgeMagnitudeSupply}.
\]
\lean{Erdos257PeriodNoncollapse.DiagonalFreshLossBridge.PowerTwoOddWindowAffine.topEdgeHalfWordBandSupply_iff_actualFinalCenteredMagnitudeSupply}{TotientActualLcmTopEdgeStaircase.lean:1479}
\quad \ev{Lean}.

\item[(iv)] $\mathrm{PowerTwoFlexibleActualTopEdgeMagnitudeSupply}$: the
same magnitude test at an \emph{arbitrary} odd rank whose adjacent depth
still fits the corridor (not forced to the canonical guarded depth):
\lean{Erdos257PeriodNoncollapse.DiagonalFreshLossBridge.PowerTwoOddWindowAffine.PowerTwoFlexibleActualTopEdgeMagnitudeSupply}{TotientActualLcmTopEdgeStaircase.lean:1927}
\ev{Open} \scale{cofinal} \coord{mobius-mersenne}. Sufficiency routes back
through (i), not through the corridor-escape chain below:
\lean{Erdos257PeriodNoncollapse.DiagonalFreshLossBridge.PowerTwoOddWindowAffine.powerTwoAdjacentSuffixMidbandSupply_of_flexibleActualTopEdgeMagnitude}{TotientActualLcmTopEdgeStaircase.lean:2007}, direct endpoint
\lean{Erdos257PeriodNoncollapse.DiagonalFreshLossBridge.PowerTwoOddWindowAffine.irrational_totientSeries_of_flexibleActualTopEdgeMagnitudeSupply}{TotientActualLcmTopEdgeStaircase.lean:2213}.

\item[(v)] $\mathrm{PowerTwoFlexibleActualTerminalDominanceSupply}$: a
strictly \emph{one-sided} version of (iv) — only the upper comparison, no
absolute value:
\[
  \mathrm{diagonalWindowIncrement}(2^a,\,2q{+}2) \le 2\cdot\mathrm{actualOddHalfCenteredLift}(a,q).
\]
\lean{Erdos257PeriodNoncollapse.DiagonalFreshLossBridge.PowerTwoOddWindowAffine.PowerTwoFlexibleActualTerminalDominanceSupply}{TotientActualLcmTopEdgeStaircase.lean:1908}
\ev{Open} \scale{cofinal} \coord{mobius-mersenne}. Direct endpoint
\lean{Erdos257PeriodNoncollapse.DiagonalFreshLossBridge.PowerTwoOddWindowAffine.irrational_totientSeries_of_flexibleActualTerminalDominanceSupply}{TotientActualLcmTopEdgeStaircase.lean:2221}.
\end{enumerate}

\paragraph{The weakest known link, and the exact identity pinning it.}
A sixth predicate, strictly weaker again than (v) — it is the disjunction of
(v)'s inequality with the opposite-direction escape — is the true minimum
of the whole cluster:
\[
  \mathrm{PowerTwoFlexibleActualTerminalCarryCorridorEscapeSupply} :\equiv
  \forall a_0,\ \exists a,q,\quad \cdots\ \wedge\
  \Bigl(2u \le d - B\ \vee\ d \le 2u\Bigr),
\]
where $d=\mathrm{diagonalWindowIncrement}(2^a,2q{+}2)$,
$u=\mathrm{actualOddHalfCenteredLift}(a,q)$,
$B=2H+(2q{+}1)+2$.
\lean{Erdos257PeriodNoncollapse.DiagonalFreshLossBridge.PowerTwoOddWindowAffine.PowerTwoFlexibleActualTerminalCarryCorridorEscapeSupply}{TotientActualLcmTopEdgeStaircase.lean:1895}
\ev{Open} \scale{cofinal} \coord{mobius-mersenne}. Direct endpoint
\lean{Erdos257PeriodNoncollapse.DiagonalFreshLossBridge.PowerTwoOddWindowAffine.irrational_totientSeries_of_flexibleActualTerminalCarryCorridorEscapeSupply}{TotientActualLcmTopEdgeStaircase.lean:2230}.

This form is not arbitrary: under integrality it is exactly one side of a
proved identity,
\[
  2\cdot\mathrm{actualOddHalfCenteredLift}(a,q)
  = \mathrm{diagonalWindowIncrement}(2^a,2q{+}2) - \mathrm{carryOrbit}(H,H,z,2q{+}1),
\]
\lean{Erdos257PeriodNoncollapse.DiagonalFreshLossBridge.PowerTwoOddWindowAffine.two_mul_actualOddHalfCenteredLift_eq_terminal_sub_trueCarry}{TotientActualLcmTopEdgeStaircase.lean:1678}
\quad \ev{Lean} \quad \scale{bounded} \quad \coord{mobius-mersenne} — an
\emph{equality}, not an inequality: the doubled centered-lift state is
exactly (terminal arithmetic letter) minus (true carry). Escaping either
side of this named open interval is exactly the corridor-escape supply.

Crucially, the sign machinery has already eliminated the branch this
identity most naturally supplies:
\lean{Erdos257PeriodNoncollapse.DiagonalFreshLossBridge.PowerTwoOddWindowAffine.actualLcm_trueEndpointSurvivor_neg}{TotientActualLcmOrbitSign.lean:172}
(\ev{Lean}, \scale{bounded}) proves that \emph{under integrality} the true
carry orbit is strictly \emph{positive} throughout the corridor, which
means the dominance branch (v) is the one the identity naturally produces
positive evidence \emph{against}, and the surviving open target is really
the \emph{lower}-escape branch,
$2u \le d - B$ — the dominance branch itself is stated separately as
\lean{Erdos257PeriodNoncollapse.DiagonalFreshLossBridge.PowerTwoOddWindowAffine.actualLcmTailOrbit_notMem_int_of_actualTerminalDominance}{TotientActualLcmTopEdgeStaircase.lean:1880}. This is the weakest exact
finite arithmetic normal form currently formalised in this corridor cluster,
not the programme's leading theorem interface: the direct fixed-full-block
first-harmonic gap above remains the constitutional outward socket.  The
corpus: a fully closed-form real-integer quantity whose escape from a named
interval is necessary and sufficient (mod a fit hypothesis
$2(H{+}q{+}2)\le 4^q$) for non-integrality at that rank.

\paragraph{A second, independent finite-window target: the short-arithmetic-kill supply.}
A separate reduction ties the diagonal form (the collapsing-free-parameters chain above) to this actual-orbit
lane via an exact digit formula with no cleanliness hypothesis at all,
$\mathrm{lcmRayArithmeticLetter}(t,j) = \varphi(H_t+j)-\varphi(H_t)$ for
\emph{every} offset $j$
(\lean{Erdos257PeriodNoncollapse.DiagonalFreshLossBridge.PowerTwoOddWindowAffine.lcmRayArithmeticLetter_eq_deltaTotient}{TotientActualLcmOrbitArithmetic.lean:298},
\ev{Lean}), and
$\mathrm{LcmDiagonalArithmeticKill}(t,L)\iff\mathrm{certifiedKill}(H_t,H_t,L)$
(\lean{Erdos257PeriodNoncollapse.DiagonalFreshLossBridge.PowerTwoOddWindowAffine.lcmDiagonalArithmeticKill_iff_certifiedKill}{TotientActualLcmOrbitArithmetic.lean:2045},
\ev{Lean}).

Consequently
\[
  \mathrm{PowerTwoActualLcmShortArithmeticKillSupply} :\equiv
  \forall a_0,\ \exists a,L,\quad a_0\le a\ \wedge\ L<2\cdot2^a\ \wedge\
  \mathrm{certifiedKill}(H_{2^a}, H_{2^a}, L)
\]
\lean{Erdos257PeriodNoncollapse.DiagonalFreshLossBridge.PowerTwoOddWindowAffine.PowerTwoActualLcmShortArithmeticKillSupply}{TotientActualLcmOrbitArithmetic.lean:2107}
\ev{Open} \scale{cofinal} \coord{mobius-mersenne} is \emph{literally the
diagonal obligation $P(t)$ of the collapsing-free-parameters chain above, restricted to powers of two and to a
short window $L<2\cdot2^a$}. It is proved sufficient for the actual-orbit
iff above via
\lean{Erdos257PeriodNoncollapse.DiagonalFreshLossBridge.PowerTwoOddWindowAffine.actualLcmOrbitNonintegralitySupply_of_shortArithmeticKill}{TotientActualLcmOrbitArithmetic.lean:2152},
and it is verified at exactly two exponents, $a=4$ (depth $L=23$) and
$a=6$ (depth $L=93$), each traced to a pre-existing compressed diagonal
certificate:
\lean{Erdos257PeriodNoncollapse.DiagonalFreshLossBridge.PowerTwoOddWindowAffine.actualLcmTailOrbit_four_and_six_notMem_int}{TotientActualLcmShortKill.lean:35}
\ev{Cert} \scale{fixed}.

The bounded stub
\lean{Erdos257PeriodNoncollapse.DiagonalFreshLossBridge.PowerTwoOddWindowAffine.powerTwoActualLcmShortArithmeticKillSupply_through_six}{TotientActualLcmShortKill.lean:54}
(\ev{Lean}, \scale{bounded}, proof is a single hard-coded witness
$(a,L)=(6,93)$, with no argument in $a$ at all) proves the supply only for
$a_0\le 6$; extending it to a single further exponent (say $a=7$) is a
self-contained, independently interesting target, not a re-run.

\paragraph{A Diophantine-flavoured alternative: the raw-approximant separation supply.}
Unconditionally, for every $a,q$:
\[
  \bigl|\mathrm{actualLcmTailOrbit}(a) - \mathrm{actualLcmRawApprox}(a,q)\bigr|
  < \frac{4H+2(2q{+}1)+4}{2^{2q+2}},
\]
\lean{Erdos257PeriodNoncollapse.DiagonalFreshLossBridge.PowerTwoOddWindowAffine.abs_actualLcmTailOrbit_sub_rawApprox_lt}{TotientActualLcmOrbitSeparation.lean:141}
\quad \ev{Lean} \quad \scale{uniform} \quad \coord{mobius-mersenne}, where
$\mathrm{actualLcmRawApprox}(a,q)$ is an explicit finite computable rational.
This reduces the whole analytic problem to a finite question: a cofinal
$1/32$-separation of the raw approximant from every integer,
\[
  \mathrm{PowerTwoActualLcmOrbitSeparationSupply} :\equiv
  \forall a_0,\ \exists a\ge\max(2,a_0),\ \exists q,\quad
  \mathrm{oddGuardedCanonicalAdjacentSuffixDepth}(2^a)=2q{+}1\ \wedge\
  \forall z\in\mathbb{Z},\ \tfrac{1}{32}+\mathrm{errorRadius}(a,q)\le|\mathrm{actualLcmTailOrbit}(a)-z|,
\]
\ev{Open} \scale{cofinal} \coord{mobius-mersenne}, sufficient via
\lean{Erdos257PeriodNoncollapse.DiagonalFreshLossBridge.PowerTwoOddWindowAffine.irrational_totientSeries_of_actualLcmOrbitSeparationSupply}{TotientActualLcmOrbitSeparation.lean:305}
(\ev{Lean}). Note the depth $q$ is not a free search parameter: the supply
pins exactly one admissible depth per scale $a$ via
$\mathrm{oddGuardedCanonicalAdjacentSuffixDepth}$.

\subsection{A problem-agnostic normal form: the guard-cylinder compression}

Independently of which coordinate is used, \emph{any} tail-difference
certificate — for arbitrary $h,N,L$, not specific to \#249 — compresses to
a two-bit test at logarithmic depth:
\[
  \bigl(\exists L,\ \mathrm{certifiedKill}(h,N,L)\bigr)
  \iff
  \mathrm{GuardCylinderWitness}(h,N).
\]
\lean{Erdos257PeriodNoncollapse.DiagonalFreshLossBridge.PowerTwoOddWindowAffine.exists_certifiedKill_iff_guardCylinderWitness}{TotientActualLcmTopEdgeStaircase.lean:844}
\quad \ev{Lean} \quad \scale{uniform} \quad \coord{seam-integer}, where the
witness is $\exists s,b,\ \mathrm{certifiedKill}(h,N{+}s,b{+}1) \vee
(\text{room} \wedge \mathrm{DyadicMixedGuard}(D(h,N{+}s,b{+}2),b))$ at
$b=\lfloor\log_2(N{+}h{+}L{+}2)\rfloor+1$. This statement mentions no
Mersenne or totient structure whatsoever and is stated for arbitrary
$h,N,L$; it compresses the search space for \emph{any} tail-difference
non-integrality certificate — of any depth $L$, however large — down to a
socket at a logarithmic scale plus a two-bit mixed-guard cylinder
($01$ or $10$). This is directly reusable, unchanged, for any binary-series
tail-difference problem, including \#257's own denominators.

\subsection{One open Farey problem and one failed rank shortcut}

The Farey growth law below is a genuinely independent live obligation.  The
rank statement is retained only to record a counterfactual sufficient input
and the reason that the generic compression route to it is closed; it is not a
second open problem supported by the present argument.

\paragraph{A counterfactual rank criterion.}
Unconditionally, for every level $e$, the dyadic totient-kernel family of
$2^e+1$ channels is linearly independent over $\mathbb{Q}$ (via CRT and
Dirichlet's theorem on primes in arithmetic progression):
\lean{Erdos249257.linearIndependent_canonicalTotientKernelFamily}{TotientMahlerDefect.lean:935}
\quad \ev{Lean} \quad \scale{uniform} \quad \coord{other:carry-kernel-rank}.
Consequently, if $S$ is rational there is a tempered integral binary carry
orbit $u$ with
\[
  \forall e,\quad 2^e-1 \le \operatorname{rank}_{\mathbb{Q}}\operatorname{span}\bigl(\mathrm{canonicalCarryKernelFamily}(u,e)\bigr).
\]
\lean{Erdos249257.not_irrational_totientSeries_implies_unbounded_carryRank_unconditional}{TotientCarryKernelRigidity.lean:300}
\quad \ev{Lean} \quad \scale{uniform} \quad \coord{other:carry-kernel-rank}.

An opposite inequality — a rationality-side rank \emph{upper} bound — would
of course contradict the displayed lower bound:
\[
  \exists C,\quad \forall\text{ tempered integral binary carry orbit } u
  \text{ for a rational } S,\quad
  \forall e,\quad \operatorname{rank}_{\mathbb{Q}}\operatorname{span}\bigl(\mathrm{canonicalCarryKernelFamily}(u,e)\bigr) \le C
\]
(or any bound growing slower than $2^e-1$).  This is a logically sufficient
new theorem schema, not an open proposition isolated by the preceding
mathematics: the corpus gives no
reason rationality should force it, and explicit finite-rank
shift-polynomial countermodels show that periodic denominator data alone do
not.

The most direct generic construction is also impossible: no
\texttt{CompressedAdjointCertificate} ($Q\cdot v\cdot A =
\mathrm{boundary}$ with $|\mathrm{boundary}|<Q\cdot v$ and $A\ne0$) can
exist —
\lean{Erdos249257.false_of_compressedAdjointCertificate}{TotientMahlerDefect.lean:1183}
\quad \ev{Lean} — and the full family is proved infinite-dimensional in
span:
\lean{Erdos249257.not_finiteDimensional_span_fullTotientKernel}{TotientMahlerDefect.lean:1145}
\quad \ev{Lean}.

What rationality \emph{does} buy — uniform eventual
periodicity of $u$'s dyadic sections modulo some $v>0$ — is proved to
\emph{not} promote to any finite rank bound without further arithmetic
input:
\lean{Erdos249257.not_irrational_totientSeries_implies_mod_period_and_unbounded_rank}{TotientTailCarryPeriod.lean:224}
\quad \ev{Lean}. This pins the obstacle precisely, but supplies no rank
upper bound.

\paragraph{Obligation 3 (Farey growth law).}
The corpus's strongest \emph{unconditional} statement about $S$ comes from
a Farey-gap denominator exclusion at a single fixed window $K=240$:
\[
  \forall p\in\mathbb{Q},\quad p.\mathrm{den} \le 79639646646701375323355774875831053
  \ \Longrightarrow\ S\ne p.
\]
\lean{Erdos249257.tsum_totient_div_pow_two_ne_ratCast_of_den_le_79639646646701375323355774875831053}{CertificateKernel.lean:18384}
\quad \ev{Lean} \quad \scale{bounded} \quad \coord{farey}. This bound is
\emph{sharp} at $K=240$ — the mediant $q=79639646646701375323355774875831054$
is proved to be the exact first failing denominator
(\lean{Erdos249257.gap_check_window_1_240_first_failure}{GapFareyBound.lean:225},
\ev{Lean}) — so re-running the same window buys nothing further; a new $K$
needs a freshly committed totient residue and freshly computed
continued-fraction convergents, both hard-coded numerals.

The cofinal
upgrade would be: a proved growth law $g(K)\to\infty$ such that for every
$K$ the $(N{=}1,K)$ gap check passes for all $q\le g(K)$ — equivalently, a
lower bound on the convergent denominators of the underlying totient-window
constant, uniform in $K$. \ev{Open} \scale{cofinal} \coord{farey}.

One
natural strengthening is proved to be a hard ceiling, not a route to
unboundedness: at a prime-power reduced denominator, the ``unit-gap''
refinement can rescue \emph{at most one} additional lattice point beyond
the ordinary gap certificate —
\lean{Erdos249257.reducedDenominatorCandidateCount_prime_pow_le_one}{PrimitiveWeightCertificate.lean:111}
\quad \ev{Lean}. This route is closed; growing $K$ genuinely requires new
per-window computation, not a refinement of the existing one.

\subsection{Equivalence map}

\begin{center}
\begin{longtable}{p{0.30\linewidth} p{0.20\linewidth} p{0.42\linewidth}}
\textbf{Form} & \textbf{Relation to Obligation 1} & \textbf{Site} \\
\hline
\endhead
Sep$(h,N,L)$ supply (Obl.\ 1, base form) & \textbf{equivalent} to Irrational$(S)$ & LcmConeFlatness.lean:412 \\
multiple-period supply & propositionally equivalent via Obl.\ 1; explicit directions are Obl.\ 1 $\Rightarrow$ multiple and multiple $\Rightarrow$ irrationality & CarrySurvivorExtinction.lean:502 \\
lcm-diagonal supply $P(t)$ & \textbf{equivalent} to Irrational$(S)$; single free parameter & LcmConeFlatness.lean:426 \\
lcm-cone supply & propositionally equivalent via Obl.\ 1; $P(t)$ is the cell $q{=}m{=}1$ & CertificateKernel.lean:18686 \\
cone-menu (\texttt{coneNonflatCert}) supply & weaker; half the pairwise radius & CertificateKernel.lean:18812 \\
pure non-integrality (diagonal/cone) & \textbf{equivalent} to the corresp.\ certificate form, via the completeness iff & CertificateKernel.lean:18706, :18718 \\
DTWFirstHarmonicNormGap & sufficient for Obl.\ 1 (not shown weaker/stronger) & FirstHarmonicPivot.lean:83 \\
actual-orbit nonintegrality supply & \textbf{equivalent} to Irrational$(S)$ itself & TotientActualLcmOrbitNonintegrality.lean:37 \\
short-arithmetic-kill supply & special case of $P(t)$: powers of two, short window & TotientActualLcmOrbitArithmetic.lean:2107 \\
top-edge residue-gap supply & sufficient for actual-orbit supply; one-sided & TotientActualLcmTopEdgeStaircase.lean:1325 \\
adjacent-suffix midband supply & weaker than top-edge residue-gap & TotientActualLcmTopEdgeStaircase.lean:1334 \\
odd-guard half-word band supply & weaker again; \textbf{equivalent} to final-magnitude form & TotientActualLcmTopEdgeStaircase.lean:1349 \\
actual final top-edge magnitude supply & \textbf{equivalent} to odd-guard half-word band & TotientActualLcmTopEdgeStaircase.lean:1471 \\
flexible top-edge magnitude supply & weaker than midband; feeds midband, not corridor-escape & TotientActualLcmTopEdgeStaircase.lean:1927 \\
flexible terminal-dominance supply & weaker again; one-sided & TotientActualLcmTopEdgeStaircase.lean:1908 \\
flexible terminal carry-corridor-escape supply & \textbf{weakest known} in this cluster & TotientActualLcmTopEdgeStaircase.lean:1895 \\
raw-approximant separation supply & alternative (Diophantine) sufficient form, same target & TotientActualLcmOrbitSeparation.lean:305 \\
rank-compression upper bound (Obl.\ 2) & independent obligation, different coordinate & TotientCarryKernelRigidity.lean:284 \\
Farey growth law (Obl.\ 3) & independent obligation, different coordinate & GapFareyBound.lean, CertificateKernel.lean:18056 \\
guard-cylinder witness & problem-agnostic normal form of \emph{any} certificate, not scale-comparable & TotientActualLcmTopEdgeStaircase.lean:844 \\
\end{longtable}
\end{center}

\emph{Reading the table.} \emph{Weaker} in rows not marked
equivalent describes an explicit implication between producer predicates:
the stronger hypothesis constructs the easier-looking one.
\emph{Equivalent} marks a registered \texttt{iff};
\emph{propositionally equivalent} marks two proved directions obtained by
composing the displayed implications with the base iff.  In particular,
Obligation~1 and the actual-orbit supply are logically equivalent through
$\mathrm{Irrational}(S)$, although the corpus does not claim a direct
witness-to-witness converter between their predicates.  The
short-arithmetic-kill supply is one concrete special case of Obligation~1's
diagonal form.

\subsection{What new mathematics each surviving form would need}

None of the forms above is close to complete; every one reduces the open
content to a named finite-flavoured arithmetic statement, never removes it.
Two forms above have their missing ingredient identified precisely enough
to name what field of technique would supply it, without any claim that the
technique exists or is easy to apply.

\paragraph{The exponential-sum form.} \texttt{DTWFirstHarmonicNormGap} asks
for cancellation in the first additive character of a totient-driven
discrepancy: a constant-saving bound, uniform in the period $h$ and growing
with the block size $X$ and depth $L$, on
$\sum_{N\in[X,2X)} e\bigl(D(h,N,L)/2^L\bigr)$, where $D$ is built from
consecutive differences of $\varphi$. This is a Weyl-sum-shaped statement
about additive character sums of an arithmetic function on a window; no
argument of this shape is attempted or proved anywhere in the corpus.

\paragraph{The staircase form.} The actual-orbit top-edge cluster's surviving
target, in every one of its six equivalent-or-sufficient guises, reduces
under the exact identity of the actual-orbit lane above to control of a single \emph{terminal
arithmetic letter}: the value $\mathrm{diagonalWindowIncrement}(2^a,2q{+}2)
= \varphi(2H{+}2q{+}2)-\varphi(2H)$ measured against twice a centred
carry-lift state. What is needed is not a new certificate mechanism — the
whole apparatus of arithmetic letters, centred lifts, and carry orbits is
already exact and unconditional — but a genuine size or residue statement
about this one totient value at cofinally many scales $a$, sufficient to
force it off the interval the identity names. The corpus's own finite
census (through every $t\le82$ on the diagonal form, and $a=4,6$ on the short
arithmetic-kill form) gives no asymptotic evidence either way; it is a
floor, not a trend.

No claim is made that either obstruction is close to resolution, and none
of the material above decides Erdős \#249, \#257, or any weakening of them.

\clearpage
\section*{Statements and declarations}

\paragraph{Artefact and data availability.} The linked Lean declarations,
fixed toolchain, and library manifest are published in the companion
repository. This manuscript provides navigation and exposition rather than
proof authority.

\paragraph{Authority boundary.} Kernel checking establishes that a proposition
was proved; it does not authorise the exposition, literature judgements, or any
claim that Problem 249 is solved.

\paragraph{Funding and competing interests.} This work received no external
funding. The author declares no competing interests.

\end{document}
